\subsection{子模；商模}

设 E 是一个 A-模，M 是 E 的一个子集；要使 E 上的 A-模结构在 M 上诱导出一个 A-模结构，必须且只需 M 是 E 的一个稳定子群（I，§4，第 3 号），因为若如此，则在 M 上诱导出的带算子群结构显然满足公理 (M \(_{II}\) )、(M \(_{III}\) ) 和 (M \(_{IV}\) )；带此结构的 M（或者，滥用说法，集合 M 本身）称为 E 的一个子模；典范单射 M → E 是线性映射。当 E 是向量空间时，其子模称为向量子空间（若不致引起混淆，简称子空间）。

例。(1) 在任意模 E 中，由 0 组成的集合是一个子模（零子模，常滥用记号记作 0）。

\begin{enumerate}
\def\labelenumi{(\arabic{enumi})}
\setcounter{enumi}{1}
\item
  设 \(\mathbf{A}\) 是一个环。 \(\mathbf{A}_s\) （相应地 \(\mathbf{A}_d\) ）的子模恰好就是环 \(\mathbf{A}\) 的左理想（相应地，右理想）。
\item
  设 E 是一个 A-模，x 是 E 的一个元素，而 a 是 A 的一个左理想。形如 \(\alpha x\) 的元素（其中 \(\alpha\) 遍历 a）组成的集合是 E 的一个子模，记作 ax。
\item
  在将交换群 G 视为 Z-模（第 1 号）时，G 的每个子群也都是子模。
\end{enumerate}

*(5) 设 I 是实直线 \(\mathbf{R}\) 上的一个开区间；在 I 上有定义且连续的实值函数组成的集合 C 是由 I 上有定义的所有实值函数组成的向量空间 \(\mathbf{R}^{\mathrm{I}}\) 的一个向量子空间。类似地，I 上可微函数的集合 D 是 \(\mathbf{C}\) 的一个向量子空间。*

设 E 是一个 A-模。与 E 上的模结构相容（I，§1，第 6 号）的每个等价关系都具有形式 \(x - y \in \mathbf{M}\) ，其中 M 是 E 的一个稳定子群（I，§4，第 4 号），即 E 的一个子模。可以立即验证，商群 E/M（I，§4，第 4 号）上的带算子群结构是一个 A-模结构，在此结构下，典范映射 E → E/M 是线性的；带此结构的 E/M 称为 E 关于子模 M 的商模。向量空间 E 的商模称为 E 的商向量空间（或简称商空间）。

例（6）。环 A 中的每个左理想 a 都定义了商模 \(A_{s}/a\) （左 A-模 \(A_{s}\) 的商模）；滥用记号，这个商模常记作 A/a。

设 E, F 为两个 A-模。由带算子群的一般性质（I，§4，第 5 号）（或直接由定义）可知，若 \(u: \mathbf{E} \to \mathbf{F}\) 是线性映射，则 E 的任意子模在 \(u\) 下的像是 F 的一个子模，而 F 的任意子模在 \(u\) 下的原像是

E 的一个子模。特别地，核 \(N = u^{-1}(0)\) 是 E 的一个子模，而 E 在 u 下的像 \(u(\mathbf{E})\) 是 F 的一个子模（I，§4，第 6 号，命题 7）；滥用说法， \(u(\mathbf{E})\) 称为 u 的像。商模 E/N 也称为 u 的余像，而商模 \(F/u(\mathbf{E})\) 称为 u 的余核。在 u 的典范分解（I，§4，第 5 号）

\[
u: \mathrm{E} \xrightarrow {p} \mathrm{E} / \mathrm{N} \xrightarrow {v} u (\mathrm{E}) \xrightarrow {i} \mathrm{F}\tag{11}
\]

v 是 u 的余像到 u 的像上的同构（第 2 号，命题 2）。u 是单射的必要且充分条件是其核为零；u 是满射的必要且充分条件是其余核为零。

u 的核、像、余像和余核分别记作 Ker u、Im u、Coim u、Coker u。

注记。设 M 是 A-模 E 的一个子模， \(\phi: \mathrm{E} \to \mathrm{E} / \mathrm{M}\) 是典范同态。要使 A-线性映射 \(u: \mathrm{E} \to \mathrm{F}\) 具有形式 \(v \circ \phi\) ，其中 \(v\) 是 \(\mathrm{E} / \mathrm{M}\) 到 \(\mathbf{F}\) 中的线性映射，必须且只需 \(\mathrm{M} \subset \operatorname{Ker}(u)\) ；因为，若此条件成立，则关系 \(x - y \in \mathbf{M}\) 蕴含 \(u(x) = u(y)\) ，从而 \(u\) 与该等价关系相容，并且显然映射 \(v: \mathrm{E} / \mathrm{M} \to \mathrm{F}\) （由 \(u\) 通过取商得到的映射）是线性的。

